%% P20 --- Gradient descent: from SGD to Adam
%%
%% Every number in this program that the reader cannot do in their head comes
%% from code/p20_gradient_descent.py and is pulled in with \val{}. The console
%% listing is a committed transcript written by that script.
%%
%% THE THESIS: every optimiser in common use is ONE update with a different
%% estimate of how far and in which direction, and each addition repairs a
%% NAMED failure of the one before it. As a list of algorithms it is six
%% recipes; as a sequence of repairs it is one line and five arguments.
%%
%% WHAT THIS PROGRAM IS OWED, read out of the files rather than remembered:
%%   F04  OWNS THE EXPONENTIAL MOVING AVERAGE OUTRIGHT -- the recurrence, the
%%        unrolled weights, the half-life, the bias correction, and the
%%        (1 - beta) convention question, which its review pass settled by
%%        measuring. F04 says in as many words: "Program P20 supplies the
%%        gradient and the rest of the optimiser; what it does not have to
%%        supply is the average." It also hands over the CONSEQUENCE of the
%%        convention -- a step size does not travel between the two forms.
%%   P15  OWNS THE ZIG-ZAG, measured on P10's bowl, and its rigour box hands
%%        the fix here by name and says why it works.
%%   P17  OWNS eta < 2/lambda_max and the bowl's optimal eta, rate and step
%%        count, so the baseline is GATED rather than recomputed. P17 also
%%        owns the rescaling argument, which turns out to be exactly the
%%        transformation the per-coordinate methods are invariant to.
%%   P11  owns the condition number; P19 owns convexity; P18 owns the
%%        gradient of the loss, so nothing here derives one.
%%
%% WHAT IT LEAVES ALONE, checked against tools/programs.json:
%%   minibatch noise, the linear scaling rule, clipping, accumulation -> P21
%%   the multiplier as the price of a constraint                      -> P22
%%   what the variance of an estimator is                             -> P24
%%
%% Twin: programs/pl/P20-gradient-descent.tex. Frame for frame, macro for
%% macro, number for number; tools/parity.py compares them.

\program[From SGD to Adam]{Gradient descent: from SGD to Adam}
\label{prog:P20}
\index{optimiser}
\index{gradient descent}

\begin{outcomes}
\outcome{1--7}{write the one update every optimiser here is a variant of, and
  say what each variant estimates differently}
\outcome{8--16}{say what averaging the recent steps repairs, and what the
  repair is worth as a function of the condition number}
\outcome{17--20}{read a momentum implementation, say which convention it uses,
  and say what has to travel with a step size}
\outcome{21--29}{say what dividing each coordinate by its own recent scale
  buys, and where the epsilon has to go}
\outcome{30--33}{correct an average that starts at zero, and say what a warmup
  is protecting against}
\outcome{34--37}{say where a penalty enters, and why that changes what it
  means}
\outcome{38--42}{read a schedule's parameters, and say what a schedule is
  doing to the step}
\end{outcomes}

\begin{quiz}
\item What does every optimiser in this program have in common?
  \teachesat{1--2}
  \answerto{The same update: move against an estimate of the downhill
  direction, by a step size. What differs between them is only how that
  estimate is formed.}
\item What sets the ceiling on a step size? \teachesat{3--4}
  \answerto{The steepest direction. Program~P17 derives
  $\eta < 2/\lambda_{\max}$, so one direction fixes the ceiling and every
  flatter direction then moves more slowly than it could.}
\item Momentum is an average of the recent gradients. What does averaging
  remove and what does it keep? \teachesat{8--9}
  \answerto{It removes what alternates in sign and keeps what does not. In a
  narrow valley the sideways component alternates every step and the forward
  one does not, so the average points along the valley.}
\item Why does Adam divide by the square root of a running average of the
  squared gradient? \teachesat{22--24}
  \answerto{Because that quantity has the units of a gradient, so the ratio
  is dimensionless and the step is about $\eta$ whatever the gradient's size.
  It makes a step size mean the same thing in every coordinate.}
\item Is weight decay the same as adding an $L_2$ penalty to the loss?
  \teachesat{35--36}
  \answerto{For plain descent, yes. For an adaptive method, no: the penalty
  added to the gradient is divided by the running scale along with everything
  else, so its effective strength differs from coordinate to coordinate.}
\item A scheduler is set with \code{gamma=0.1}. What is the learning rate
  after three boundaries? \teachesat{38--39}
  \answerto{A thousandth of what it started at, because it multiplies. From
  $\num{1e-3}$ it goes to $\num{1e-6}$, not to $\num{0.1}$.}
\end{quiz}

% ============================================================
\section{One update, and what has to be estimated}
\label{sec:P20-one-update}
\index{optimiser!the common update}

\begin{fr}
Programs \ref{prog:P15}, \ref{prog:P17} and \ref{prog:P19} have each left
something here. \ref{prog:P15} produced a walk that wastes most of its steps
and said so. \ref{prog:P17} derived the ceiling on a step size and said the
optimisers were this program's. \ref{prog:P19} said what is lost when the
surface is not convex, and what is not.

Before any of that, the object they are all about. Write down the update that
plain gradient descent performs, and then say what a reader has to choose
before it can run.
\dotline
\nextframe
\end{fr}

\begin{fr}
\begin{ansblock}
\[ w \;\leftarrow\; w - \eta\,\nabla f(w) \]
and the choice is $\eta$ \dash{} how far.
\end{ansblock}

Now the observation the whole program rests on. Every optimiser in common use
writes that same line. What changes between them is not the shape of the
update but \textbf{what stands in the place of $\nabla f(w)$}: the last
gradient, an average of the recent ones, or one divided by its own recent
size.

\[ w \;\leftarrow\; w - \eta\,d, \qquad d = \text{an estimate of where to go} \]

\mermaidfig{p20-one-update}{One line, and what varies between the methods.}{one
line, several estimates}
\end{fr}

\begin{fr}
That is worth taking seriously rather than treating as a slogan, because it
changes what there is to learn. A list of six algorithms is six things to
memorise. One line with five repairs is one thing, plus five arguments you can
reconstruct.

So the question for the rest of this program is not \emph{what does Adam do}
but \emph{what was wrong with the thing before it}. Start with what is wrong
with the line as it stands. Program~\ref{prog:P17} derived a bound on $\eta$:
what was it, and what fixes it?
\dotline
\nextframe
\end{fr}

\begin{fr}
\begin{ansblock}
$\eta < 2/\lambda_{\max}$ \dash{} and it is fixed by the \emph{steepest}
direction.
\end{ansblock}

Which is the whole problem in one line. The ceiling is set by the direction
that curves most, and every other direction has to live under that ceiling. A
direction with curvature $\lambda$ moves by a factor of $1 - \eta\lambda$ each
step, so the flattest one creeps while the steepest one is at the edge of
stability.

On the bowl Programs \ref{prog:P10}, \ref{prog:P15} and \ref{prog:P17} all
use \dash{} curvatures $\val{p20.lam.hi}$ and $\val{p20.lam.lo}$, a condition
number of $\val{p20.kappa}$ \dash{} the best step size is
$\val{p20.eta.best}$, the walk shrinks by $\val{p20.rate.sgd}$ each step, and
it needs $\val{p20.steps.sgd}$ steps to close $99$ per cent of the distance.

\begin{note}
Those three numbers are Program~\ref{prog:P17}'s, and this program's script
reads them out of that program's committed values and fails the build if they
disagree. Four programs are now talking about one bowl rather than four that
resemble each other.
\end{note}
\end{fr}

\begin{fr}
$\val{p20.steps.sgd}$ steps to cover a distance of about $\num{1.4}$. Where is
the effort going? Program~\ref{prog:P15} measured it: over the first
$\val{p15.zig.steps}$ steps the steep coordinate changes sign
$\val{p15.zig.crossings}$ times, and the walk travels $\val{p15.zig.ratio}$
times as far as it moves.

So say what a step of that walk consists of, in two parts.
\nextframe
\end{fr}

\begin{fr}
\begin{ansblock}
A large sideways component that reverses every step, and a small forward one
that does not.
\end{ansblock}

The sideways part is not merely wasted. It is \emph{consistently} wasted: it
points one way, then the other, then back again, and over two steps it very
nearly cancels. Something that nearly cancels over two steps is exactly the
kind of thing an average removes, which is the next section and is
Program~\ref{prog:P15}'s own suggestion.

\begin{trapbox}
\textbf{\enquote{The gradient is large, so I am far from the minimum.}} On
this bowl the gradient is large and most of it is pointing across the valley
rather than along it. Program~\ref{prog:P15} says this in full; it is repeated
here because it is the reason a step count and a gradient norm disagree, and
because every repair in this program is a repair to the \emph{direction}
rather than to the distance.
\end{trapbox}
\end{fr}

\begin{fr}
One more thing to fix in place before the repairs start, because it is the
thing that makes them comparable.

For a quadratic, plain descent at its best step size has a known convergence
rate: each step multiplies the distance by
\[ \frac{\kappa - 1}{\kappa + 1} \]
where $\kappa$ is the condition number Program~\ref{prog:P11} defined. At
$\kappa = \val{p20.kappa}$ that is $\val{p20.rate.sgd}$, so reaching one per
cent of the starting distance needs $\val{p20.pred.sgd}$ steps.

The measurement is $\val{p20.steps.sgd}$. \textbf{The prediction is exact}, and
the script checks it at $\val{p20.sweep.kappas}$ condition numbers from $\val{p20.sweep.k.lo}$ to
$\val{p20.sweep.k.hi}$ rather than at this one. Hold on to that, because the
next section's prediction behaves differently and the difference is
instructive.
\end{fr}

% ============================================================
\section{Averaging the direction}
\label{sec:P20-momentum}
\index{optimiser!momentum}
\index{moving average!as momentum}

\begin{fr}
Program~\ref{prog:F04} built an exponential moving average out of arithmetic
on a list of numbers, took it apart completely, and closed by saying that
\textbf{momentum is that average with a gradient in it}.

So there is nothing to derive here. Write the recurrence, with $g$ for the
gradient, and say what you would step along.
\dotline
\nextframe
\end{fr}

\begin{fr}
\begin{ansblock}
\[ v \;\leftarrow\; \beta v + (1-\beta) g, \qquad
   w \;\leftarrow\; w - \eta v \]
\end{ansblock}

Step along the average rather than along the last gradient. And now the
argument Program~\ref{prog:P15} left standing: over the recent steps the
sideways component alternates in sign and the forward one does not, so the
average keeps the forward one and cancels the sideways one against itself.

\mermaidfig{p20-each-fixes-one}{Each addition repairs the one before it.}{a
sequence of repairs}
\end{fr}

\begin{fr}
That is a mechanism, and a mechanism is not a number. For a quadratic the
right coefficients are known in closed form \dash{} on this bowl
$\beta = \val{p20.beta.best}$ and $\eta = \val{p20.eta.mom}$ \dash{} and the
rate they give is
\[ \frac{\sqrt{\kappa} - 1}{\sqrt{\kappa} + 1} \]
against plain descent's $(\kappa-1)/(\kappa+1)$.

Compare the two expressions and say, in one word, what momentum buys.
\nextframe
\end{fr}

\begin{fr}
\begin{ansblock}
A \textbf{square root}. Momentum replaces $\kappa$ by $\sqrt{\kappa}$.
\end{ansblock}

Which is a statement about the \emph{problem} rather than about any
implementation, and that is why it can be predicted before it is measured. At
$\kappa = \val{p20.kappa}$ the rate falls from $\val{p20.rate.sgd}$ to
$\val{p20.rate.mom}$, so the predicted step count falls from
$\val{p20.pred.sgd}$ to $\val{p20.pred.mom}$.

The measurement is $\val{p20.steps.mom}$.
\end{fr}

\begin{fr}
$\val{p20.steps.mom}$ against a prediction of $\val{p20.pred.mom}$. That is a
long way off, and the draft of this program asserted the two would agree.

Before reading on: the prediction is not wrong and the measurement is not
wrong. What could a \emph{rate} fail to describe about a step count?
\dotline
\nextframe
\end{fr}

\begin{fr}
\begin{ansblock}
The beginning. A rate says what each step does to the one before it once the
walk has settled; it says nothing about how the walk starts.
\end{ansblock}

And this walk starts by going the wrong way. Measured, the distance to the
minimum \emph{rises} to $\val{p20.overshoot}$ times its starting value \dash{}
$\val{p20.overshoot.pct}$ per cent further away \dash{} before it comes down.
The average has to be built out of gradients before it can be useful, and while
it is being built it carries the walk past.

Once it has settled, the rate is right: measured over fifty steps of the tail,
the distance falls by $\val{p20.tail.rate}$ each step against a predicted
$\val{p20.rate.mom}$. So on this bowl momentum is worth
$\val{p20.speedup}$ times plain descent, against a $\sqrt{\kappa}$ of
$\val{p20.sqrt.kappa}$.
\end{fr}

\begin{fr}
\begin{rigourbox}
There is a second reason, and it is exact rather than a transient. The
optimal coefficients make the two roots of the iteration \emph{coincide} in
both eigendirections at once \dash{} the discriminant is
$\val{p20.disc}$, which the script checks \dash{} and that is precisely what
makes them optimal. A repeated root with a single eigenvector decays like
$k\rho^{k}$ rather than $\rho^{k}$, so the predicted rate is approached from
above and never quite attained.

The algebra is a first course's, and this program does not do it. What is
worth keeping is the shape of the conclusion: \textbf{a rate is a limit, and
a step count is not}.
\end{rigourbox}
\end{fr}

\begin{fr}
So the honest comparison is the one swept over several problems rather than
quoted on one. Over condition numbers from $\val{p20.sweep.k.lo}$ to
$\val{p20.sweep.k.hi}$:

\begin{center}
\begin{tabular}{lll}
\toprule
& plain descent & momentum \\
\midrule
prediction & exact at every $\kappa$ & a floor it approaches \\
$\kappa = \val{p20.sweep.k.lo}$ & \dash{} & $\val{p20.sweep.ratio.lo}\times$ \\
$\kappa = \val{p20.sweep.k.hi}$ & \dash{} & $\val{p20.sweep.ratio.hi}\times$ \\
ceiling & \dash{} & $\sqrt{\kappa}$, here $\val{p20.sweep.sqrt.hi}$ \\
\bottomrule
\end{tabular}
\end{center}

The advantage grows with the condition number and never reaches its ceiling.
Take the first half of that: on which problems is momentum barely worth
having, and why?
\nextframe
\end{fr}

\begin{fr}
\begin{ansblock}
Well conditioned ones. At $\kappa = \val{p20.sweep.k.lo}$ the advantage is
$\val{p20.sweep.ratio.lo}$, because $\sqrt{\kappa}$ is close to $\kappa$ when
$\kappa$ is small.
\end{ansblock}

The second half is the shortfall, and it is why $\sqrt{\kappa}$ is a bound
rather than a promise.

\begin{trapbox}
\textbf{\enquote{Momentum gives a $\sqrt{\kappa}$ speedup.}} It gives a rate
that improves like $\sqrt{\kappa}$, which is not the same claim. Measured
here, the step count improves by $\val{p20.sweep.ratio.hi}$ at
$\kappa = \val{p20.sweep.k.hi}$ against a $\sqrt{\kappa}$ of
$\val{p20.sweep.sqrt.hi}$ \dash{} less than half of it.

The gap is the transient, and it does not go away by tuning: it is what a
repeated root costs. Quote the rate, or quote a measured count. Quoting the
rate as though it were a count is how a bound becomes folklore.
\end{trapbox}
\end{fr}

% ============================================================
\section{The convention that changes the step}
\label{sec:P20-convention}
\index{optimiser!momentum conventions}

\begin{fr}
Program~\ref{prog:F04} flagged this and handed the consequence here. Two
implementations of momentum:

\begin{center}
\begin{tabular}{ll}
\toprule
A & $v \leftarrow \beta v + (1-\beta) g$ \\
B & $v \leftarrow \beta v + g$ \\
\bottomrule
\end{tabular}
\end{center}

Both are called momentum and both appear in real code. Are they the same
update?
\dotline
\nextframe
\end{fr}

\begin{fr}
\begin{ansblock}
The same \emph{direction}, and not the same \emph{length}. B's vector is
$1/(1-\beta)$ times A's.
\end{ansblock}

At $\beta = \val{p20.conv.beta}$ that factor is $\val{p20.fold}$. Since a
constant multiplying the step is indistinguishable from a change of step size,
neither form is wrong and both are in use. What does not survive is a step size
carried from one to the other.

Measured on the same bowl at $\eta = \val{p20.conv.eta}$ and
$\beta = \val{p20.conv.beta}$: form A converges in $\val{p20.conv.same}$ steps,
form B with $\eta$ divided by $\val{p20.fold}$ reproduces that walk
\textbf{exactly}, step for step \dash{} and form B with the same $\eta$
diverges.
\end{fr}

\begin{fr}
\begin{warning}
\textbf{Not \enquote{slower}. Gone.} Ten times a working step size on this
bowl is past Program~\ref{prog:P17}'s bound, so the walk leaves the basin and
never returns.

Nothing about the name tells you which form you have, and nothing about the
symbol does either. \textbf{Read the line that updates the state}: which side
the coefficient sits on, and whether the other side carries a factor at all.
Program~\ref{prog:F04} makes the same point about which value $\beta$ weights;
this is the second half of it and it is the half that costs a run.
\end{warning}

So when a hyperparameter is carried from one library to another, what has to
travel with it?
\nextframe
\end{fr}

\begin{fr}
\begin{ansblock}
The expression it appears in, not the number.
\end{ansblock}

A hyperparameter is only meaningful together with the line it sits in. Two libraries can both offer \code{momentum=0.9} and mean different
lengths; two schedulers can both offer \code{gamma} and mean different
operations, which is section 7. When a number is carried from one setting to
another, what has to travel with it is not the number but the line it sits in.
\end{fr}

% ============================================================
\section{Scaling each coordinate}
\label{sec:P20-adaptive}
\index{optimiser!per-coordinate scaling}
\index{optimiser!Adam}

\begin{fr}
Momentum repaired the direction. The step size is still one number for every
coordinate at once, and that is the next failure.

Suppose two parameters whose gradients are routinely a hundred times apart in
size \dash{} an embedding row touched by every batch and one touched by a rare
token, say. A single $\eta$ has to serve both. What is the difficulty, stated
precisely?
\dotline
\nextframe
\end{fr}

\begin{fr}
\begin{ansblock}
A step size that is right for one is wrong for the other by the same factor of
a hundred: too small to make progress, or too large to be stable.
\end{ansblock}

So the repair has to make a step size \emph{mean the same thing} in every
coordinate. Which is a question about units. The update $\eta g$ has the units
of $\eta$ times a gradient, so its size depends on the gradient's size, and
that is exactly what we do not want.

What could you divide a gradient by so that what is left has no units at all?
\nextframe
\end{fr}

\begin{fr}
\begin{ansblock}
Something that is itself a gradient. The square root of an average of the
squared gradient is one.
\end{ansblock}

\[ v \leftarrow \beta_2 v + (1-\beta_2) g^{2}, \qquad
   w \leftarrow w - \eta\,\frac{g}{\sqrt{v} + \varepsilon} \]

$\sqrt{v}$ has the units of a gradient, so $g/\sqrt{v}$ has none, and
\textbf{the step is about $\eta$ however large the gradient is}.

Now put section 2's average on the numerator as well, so that both the
direction and the scale are running averages. Name what you have written.
\dotline
\nextframe
\end{fr}

\begin{fr}
\ans{Adam}
The same recurrence run twice \dash{} once on the gradient and once on its
square \dash{} which is exactly what Program~\ref{prog:F04} said it would be,
two parts before any of the calculus arrived.

That claim is checkable, so it is checked. Starting from a cold state, the
first step is taken for gradients spanning $\val{p20.adam.decades}$ decades:

\transcript{p20-unit-step}

The step is $\val{p20.adam.eta}$ in every coordinate for the large gradients,
and short of it by $\val{p20.adam.shortfall}$ per cent for the small ones.
\textbf{This is why Adam is insensitive to loss scaling}: multiply the loss by
a thousand and every gradient is multiplied by a thousand, and the numerator
and the denominator are multiplied by a thousand together.
\end{fr}

\begin{fr}
The shortfall is not noise and it is not an approximation. The first step is
exactly
\[ \frac{\eta}{1 + \varepsilon/\lvert g\rvert} \]
which the script checks in every coordinate at every scale. So the departure
from $\eta$ is entirely the epsilon's, and it is largest where the gradient is
\emph{smallest} \dash{} the opposite end from the one people worry about.

Which raises the question the epsilon is usually written without any thought
about. It is there to stop a division by zero. Where should it go: inside the
square root, added to $v$, or outside it, added to $\sqrt{v}$?
\dotline
\nextframe
\end{fr}

\begin{fr}
\begin{ansblock}
Outside. Inside, it is added to a gradient \emph{squared}, which is not a
gradient.
\end{ansblock}

That is the units argument again and it has a measurable consequence. With
$\varepsilon = \val{p20.eps}$ and a gradient of $\val{p20.eps.g}$:

\begin{center}
\begin{tabular}{ll}
\toprule
$\varepsilon$ outside the root & step short by $\val{p20.eps.short.out}$ per cent \\
$\varepsilon$ inside the root & step short by $\val{p20.eps.short.in}$ per cent \\
\bottomrule
\end{tabular}
\end{center}

Outside, the epsilon is compared with $\lvert g \rvert$ and its effect is
$\varepsilon/\lvert g\rvert$. Inside, it is compared with $g^{2}$ and its
effect is of order $\varepsilon/g^{2}$, which is not small when the gradient is
small \dash{} and a small-gradient coordinate is exactly the one the whole
mechanism exists to rescue.
\end{fr}

\begin{fr}
\begin{trapbox}
\textbf{\enquote{Epsilon is a numerical guard, so its exact placement does not
matter.}} It decides what the guard is being compared against. Written inside
the root it silently becomes a floor on the gradient's \emph{square}, so the
coordinates it damps are chosen by the loss's units rather than by anything
about the model \dash{} and it damps them by
$\val{p20.eps.short.in}$ per cent where the correct placement costs
$\val{p20.eps.short.out}$.

A guard that changes the answer is not a guard.
\end{trapbox}
\end{fr}

\begin{fr}
There is a sharper way to say what the division buys, and it uses something
Program~\ref{prog:P17} has already established.

That program showed that writing a parameter as $w = c\,u$ leaves the function
the network computes completely unchanged, while multiplying that direction's
curvature by $c^{2}$. Same model, different coordinates.

Take $c = \val{p20.rescale.c}$ on the steep direction of our bowl. What happens
to plain descent at the step size that was optimal before?
\nextframe
\end{fr}

\begin{fr}
\begin{ansblock}
It diverges. The curvature is multiplied by $\val{p20.rescale.csq}$, so the
step size is now $\val{p20.sgd.overshoot}$ times past
Program~\ref{prog:P17}'s bound.
\end{ansblock}

And Adam takes $\val{p20.adam.plain}$ steps before the reparameterisation and
$\val{p20.adam.rescaled}$ after it \dash{} the same number, because both the
gradient and its running scale are multiplied by $c$ and the ratio is
untouched.

\textbf{Plain descent's answer depends on how somebody chose to write the
parameters. Adam's does not.} That is a stronger statement than a speed
comparison, because it does not depend on the problem.
\end{fr}

% ============================================================
\section{Correcting an average that starts at zero}
\label{sec:P20-bias}
\index{moving average!bias correction}
\index{optimiser!warmup}

\begin{fr}
Both of those averages start at zero, and Program~\ref{prog:F04} took that
apart too. Its result, in one line: after $t$ steps an average of a constant
reads $1 - \beta^{\,t}$ of the truth, so dividing by $1 - \beta^{\,t}$ puts it
back.

Adam uses two coefficients, $\beta_1 = \num{0.9}$ on the gradient and
$\beta_2 = \num{0.999}$ on its square. Their half-lives are
$\val{f04.ema.halflife.steps}$ and $\val{f04.ema.b999.steps}$ steps. What does
that say about the correction?
\dotline
\nextframe
\end{fr}

\begin{fr}
\begin{ansblock}
That the two moments are biased by very different amounts, so each needs its
own $1 - \beta^{\,t}$ \dash{} and that the second one is doing real work for
hundreds of steps rather than for three.
\end{ansblock}

Which is the whole of Adam's bias correction, and it arrived without any
calculus, in a Foundation program, out of a geometric sum.

\begin{note}
Notice what the correction is \emph{not}. It is not a fudge for a slow start
and it is not tuned. It is the exact reciprocal of the fraction of the weight
that has arrived, so on a constant sequence it is right at every $t$ in the
algebra \dash{} and right to about the last place in binary64, which
Program~\ref{prog:F04} measured.
\end{note}
\end{fr}

\begin{fr}
The correction fixes the average's \emph{scale} early on. It does not fix its
\emph{quality}, and the difference is what the last piece of Adam's usual
recipe is about.

After one step, $v$ has seen one gradient. After ten, ten. The denominator is
an estimate of a coordinate's typical gradient size built from that sample.
Say what is wrong with it early.
\nextframe
\end{fr}

\begin{fr}
\begin{ansblock}
It is an estimate from almost no data, so it moves a great deal from step to
step \dash{} and it is in a denominator.
\end{ansblock}

Measured on a stream whose gradient sizes vary by a factor of $25$: over the
first ten steps the corrected estimate swings by a factor of
$\val{p20.warm.early}$, and after three hundred steps by a factor of
$\val{p20.warm.late}$.

\textbf{That is what a warmup is protecting against.} Not a fragile model and
not a bad initialisation \dash{} a denominator estimated from a handful of
samples, being used to divide by. Raising $\eta$ over the first few hundred
steps keeps the step small exactly while the quantity it is divided by is
least trustworthy.
\end{fr}

% ============================================================
\section{Where the penalty enters}
\label{sec:P20-decay}
\index{optimiser!weight decay}
\index{optimiser!AdamW}

\begin{fr}
Two more repairs, and both are about things that sit \emph{around} the update
rather than inside it. The first is the penalty.

There are two places a weight penalty can be applied: added to the gradient
before the optimiser sees it, or subtracted from the weight after the
optimiser has finished. Write both, and say whether they differ for plain
gradient descent.
\dotline
\nextframe
\end{fr}

\begin{fr}
\begin{ansblock}
\[ w \leftarrow w - \eta(g + \lambda w) \qquad\text{against}\qquad
   w \leftarrow w - \eta g - \eta\lambda w \]
and for plain descent they are the same line rearranged.
\end{ansblock}

Which is why the two names have been used interchangeably for years. The
script settles it: with plain descent both settle in the same place to better
than $\val{p20.wd.sgd.gap}$.

\mermaidfig{p20-penalty-enters}{Two places to apply the same penalty.}{where a
penalty enters}

Now put an adaptive method in between them. What happens to the first form?
\nextframe
\end{fr}

\begin{fr}
\begin{ansblock}
The $\lambda w$ term is divided by $\sqrt{v}$ along with everything else, so
its effective strength becomes $\lambda/\sqrt{v}$ \dash{} a different number in
every coordinate, and a moving one.
\end{ansblock}

Measured. Two coordinates differing only in curvature, by a factor of
$\val{p20.wd.curvratio}$, both with their minimum at $1$, one penalty
$\lambda = \val{p20.wd.lambda}$:

\begin{center}
\begin{tabular}{lll}
\toprule
& steep coordinate & flat coordinate \\
\midrule
penalty in the gradient & $\val{p20.wd.l2.steep}$ & $\val{p20.wd.l2.flat}$ \\
penalty on the weight & $\val{p20.wd.decoupled}$ & $\val{p20.wd.decoupled}$ \\
\bottomrule
\end{tabular}
\end{center}

The same $\lambda$ shrinks one coordinate by nine per cent and the other by
ninety, a factor of $\val{p20.wd.l2.spread}$ \dash{} and the one it nearly
annihilates is the one with the small gradient. The decoupled row is two
identical numbers to better than $\val{p20.wd.gap}$, because that penalty
knows nothing about either coordinate.
\end{fr}

\begin{fr}
\begin{trapbox}
\textbf{\enquote{Adam with \code{weight\_decay} is $L_2$ regularisation.}} It
is not, and the table above is why. Adding the penalty to the gradient makes
its effective strength a function of each parameter's own gradient history, so
\enquote{the regularisation strength} is not one number and does not stay
still.

That is the entire content of the distinction between the two \dash{} one
line moved from before the optimiser to after it \dash{} and the reason a
decay coefficient tuned with one optimiser does not transfer to the other.
\end{trapbox}
\end{fr}

% ============================================================
\section{What a schedule does to the step}
\label{sec:P20-schedules}
\index{optimiser!learning-rate schedule}

\begin{fr}
The last piece is the one usually described as ritual: the step size is not
held fixed but driven along a curve.

Take the commonest form first, and take it as a reading exercise rather than
as a design question. A scheduler is configured with a starting rate of
$\val{p20.sched.lr0}$ and \code{gamma=\val{p20.sched.gamma}}, and the run
passes $\val{p20.sched.boundaries}$ boundaries. What is the learning rate at
the end?
\dotline
\nextframe
\end{fr}

\begin{fr}
\begin{ansblock}
$\val{p20.sched.after}$.
\end{ansblock}

\begin{trapbox}
\textbf{\enquote{\code{gamma=0.1} means the learning rate goes to
$\num{0.1}$.}} It \emph{multiplies}. Three boundaries take
$\val{p20.sched.lr0}$ to $\val{p20.sched.after}$, which is
$\val{p20.sched.factor}$ times smaller than the misreading, and the two
answers are five orders of magnitude apart.

That is why the symptom is \enquote{the model stopped learning at epoch
$30$} rather than \enquote{it learned a little more slowly}: the step size did
not settle at a floor, it kept going. And no two schedule families are
parameterised alike, so the reading has to be done again for each one.
\end{trapbox}
\end{fr}

\begin{fr}
Two schedules, over the same run length: a cosine from the peak to zero, and a
straight line from the peak to zero. Before reading on \dash{} which spends
more of the total budget, and how much more?
\nextframe
\end{fr}

\begin{fr}
\begin{ansblock}
Neither. They spend the same, to better than $\val{p20.cos.lin.gap}$ of the
peak.
\end{ansblock}

Both average $\val{p20.cos.area}$ of the peak rate over the run, because a
cosine from $1$ to $0$ is symmetric about its midpoint and so is a line. What
differs is entirely \emph{where} the budget is spent: the cosine is still at
$\val{p20.cos.quarter}$ of the peak a quarter of the way through, where the
line is at $\num{0.75}$, and it gives that back at the end.

So a schedule is not a way of using less step. It is a way of choosing when to
use it \dash{} large while the surface is being explored, small while it is
being settled into \dash{} and the argument for one shape over another has to
be made in those terms or not at all.
\end{fr}

\begin{fr}
Which closes the program, and it is worth saying what has \emph{not} been
claimed.

Every repair here was derived against a named failure and, where it could be,
measured on one bowl. None of it says that Adam is better than momentum, or
that a cosine is better than a line, on the surfaces this book's readers
actually train on. Those are empirical questions about a class of problem that
nobody has characterised, and Program~\ref{prog:P19} has already said why the
theory that would settle them does not apply.

What the program does leave you with is the ability to read an optimiser's
update line and say which failure each term is there for \dash{} and to notice
when a hyperparameter has been carried across a change that alters what it
means.
\end{fr}

\begin{summarybox}
\sumitem{1--2}{Every optimiser here is $w \leftarrow w - \eta d$ with a
  different $d$. \textbf{The list is one line and a sequence of repairs},
  not six algorithms.}
\sumitem{3--4}{The step size is capped by the \emph{steepest} direction,
  $\eta < 2/\lambda_{\max}$, so every flatter direction moves more slowly
  than it could. That is Program~\ref{prog:P17}'s bound, and it is what
  everything after it repairs.}
\sumitem{5--6}{The wasted effort is a sideways component that reverses every
  step: $\val{p15.zig.crossings}$ sign changes in $\val{p15.zig.steps}$ steps,
  which is Program~\ref{prog:P15}'s measurement on the same bowl.}
\sumitem{7}{On that bowl plain descent needs $\val{p20.steps.sgd}$ steps, and
  its predicted count from $(\kappa-1)/(\kappa+1)$ is \textbf{exact at every
  condition number tried}.}
\sumitem{8--9}{Momentum is Program~\ref{prog:F04}'s average with a gradient
  in it: it cancels what alternates and keeps what does not, which is exactly
  the shape of the zig-zag.}
\sumitem{10--11}{It replaces $\kappa$ by $\sqrt{\kappa}$ \textbf{in the
  rate} \dash{} $\val{p20.rate.sgd}$ becomes $\val{p20.rate.mom}$.}
\sumitem{12--14}{The measured count is $\val{p20.steps.mom}$ against a
  predicted $\val{p20.pred.mom}$, because the walk first overshoots to
  $\val{p20.overshoot}$ times its starting distance and because the optimal
  coefficients give a repeated root. \textbf{A rate is a limit and a step
  count is not.}}
\sumitem{15--16}{The advantage grows with $\kappa$ and never reaches
  $\sqrt{\kappa}$: $\val{p20.sweep.ratio.lo}\times$ at
  $\kappa = \val{p20.sweep.k.lo}$ and $\val{p20.sweep.ratio.hi}\times$ at
  $\kappa = \val{p20.sweep.k.hi}$, against a ceiling of
  $\val{p20.sweep.sqrt.hi}$.}
\sumitem{17--19}{$v \leftarrow \beta v + g$ and
  $v \leftarrow \beta v + (1-\beta)g$ point the same way and differ in length
  by $1/(1-\beta)$. \textbf{A step size does not travel between them}:
  rescaling reproduces a $\val{p20.conv.same}$-step walk exactly and not
  rescaling diverges.}
\sumitem{20}{What travels with a hyperparameter is the expression it appears
  in, never the number on its own.}
\sumitem{21--24}{Dividing by $\sqrt{v}$ removes the units, so the step is
  about $\eta$ whatever the gradient's size \dash{} which is also why Adam is
  insensitive to loss scaling. Both averages together are Adam.}
\sumitem{25}{Measured over twelve decades the first step is
  $\eta/(1 + \varepsilon/\lvert g\rvert)$, so the only departure from $\eta$
  is the epsilon's, and it is largest where the gradient is smallest.}
\sumitem{26--27}{The epsilon goes \textbf{outside} the root, where it is
  compared with a gradient. Inside it is compared with a gradient squared:
  $\val{p20.eps.short.in}$ per cent of the step lost against
  $\val{p20.eps.short.out}$.}
\sumitem{28--29}{Under Program~\ref{prog:P17}'s reparameterisation plain
  descent diverges \dash{} $\val{p20.sgd.overshoot}$ times past its bound
  \dash{} and Adam takes the same $\val{p20.adam.plain}$ steps.
  \textbf{An adaptive method's answer does not depend on how the parameters
  were written.}}
\sumitem{30--31}{Bias correction is Program~\ref{prog:F04}'s
  $1 - \beta^{\,t}$, one per coefficient, and at $\beta_2 = \num{0.999}$ it is
  working for hundreds of steps rather than for three.}
\sumitem{32--33}{A warmup protects a denominator estimated from almost no
  data: the scale estimate swings by $\val{p20.warm.early}$ over ten steps and
  $\val{p20.warm.late}$ over three hundred.}
\sumitem{34--35}{A penalty can be added to the gradient or subtracted from
  the weight, and \textbf{for plain descent the two are the same line}.}
\sumitem{36--37}{For an adaptive method they are not: the penalty in the
  gradient is divided by the running scale with everything else, so one
  $\lambda$ leaves the steep coordinate at $\val{p20.wd.l2.steep}$ and the
  flat one at $\val{p20.wd.l2.flat}$.}
\sumitem{38--39}{\code{gamma} \emph{multiplies}:
  $\val{p20.sched.lr0}$ becomes $\val{p20.sched.after}$ over
  $\val{p20.sched.boundaries}$ boundaries, which is
  $\val{p20.sched.factor}$ times smaller than the misreading.}
\sumitem{40--42}{A cosine and a line spend the same budget \dash{}
  $\val{p20.cos.area}$ of the peak \dash{} so the argument for a shape is
  about \emph{when} the step is spent, never about \emph{how much}.}
\end{summarybox}

\canyou

\begin{testexercises}
\item A quadratic has curvatures $8$ and $2$. What is the largest step size
  plain gradient descent can use, and what is the condition number?
  \answerto{The bound is $2/\lambda_{\max} = \num{0.25}$, and the condition
  number is $8/2 = 4$.}
\item An implementation reads \code{v = 0.9*v + grad} and is being ported to
  one that reads \code{v = 0.9*v + 0.1*grad}. The old step size was
  $\num{0.001}$. What should the new one be, and why?
  \answerto{$\num{0.01}$. The second form's velocity is $1/(1-\beta) = 10$
  times shorter, so the step size has to be ten times larger to produce the
  same walk. Carrying $\num{0.001}$ across would make every step a tenth of
  its intended length.}
\item Adam is running on a loss that is then multiplied by $1000$. What
  happens to the size of its steps, and why?
  \answerto{Nothing, to the epsilon. Every gradient is multiplied by $1000$,
  so the numerator and $\sqrt{v}$ in the denominator are both multiplied by
  $1000$ and the ratio is unchanged. It is the property that makes a step size
  transferable between losses of different scales.}
\item Two coordinates have gradients that are routinely a factor of $100$
  apart. An $L_2$ penalty of $\lambda$ is added to the loss and the optimiser
  is Adam. Which coordinate is regularised harder, and by roughly how much?
  \answerto{The one with the \emph{smaller} gradient, by about the same factor
  of $100$: the penalty passes through the $1/\sqrt{v}$ division, so its
  effective strength is $\lambda/\sqrt{v}$, which is largest where $\sqrt{v}$
  is smallest.}
\item A schedule starts at $\num{2e-4}$ with \code{gamma=0.5} and passes four
  boundaries. What is the final rate?
  \answerto{$\num{2e-4} \times \num{0.5}^{4} = \num{1.25e-5}$. It multiplies
  once per boundary; it does not approach $\num{0.5}$.}
\end{testexercises}

\begin{furtherproblems}
\item Show that for a quadratic with curvatures $\lambda_{\min}$ and
  $\lambda_{\max}$ the step size $2/(\lambda_{\min} + \lambda_{\max})$ makes
  the two extreme directions shrink by the same factor, and that this factor
  is $(\kappa - 1)/(\kappa + 1)$.
  \answerto{A direction with curvature $\lambda$ is multiplied by
  $1 - \eta\lambda$ each step. At $\eta = 2/(\lambda_{\min}+\lambda_{\max})$
  the two factors are
  \[ 1 - \frac{2\lambda_{\min}}{\lambda_{\min}+\lambda_{\max}}
     \quad\text{and}\quad
     1 - \frac{2\lambda_{\max}}{\lambda_{\min}+\lambda_{\max}} \]
  which are
  $\pm(\lambda_{\max}-\lambda_{\min})/(\lambda_{\max}+\lambda_{\min})$
  \dash{} equal in size and opposite in sign. Dividing top and bottom by
  $\lambda_{\min}$ gives $(\kappa-1)/(\kappa+1)$. Any smaller $\eta$ leaves
  the flat direction slower; any larger one makes the steep factor bigger in
  size.}
\item Adam's update is $\eta\,\hat m/(\sqrt{\hat v} + \varepsilon)$. Show that
  multiplying every gradient by a constant $c > 0$ leaves the update
  unchanged in the limit of small $\varepsilon$, and say what breaks the
  invariance when $\varepsilon$ is not small.
  \answerto{$\hat m$ is linear in the gradients, so it scales by $c$;
  $\hat v$ is quadratic, so $\sqrt{\hat v}$ scales by $c$; the ratio is
  unchanged. The epsilon is the only term that does not scale, so the update
  becomes $\eta c\hat m/(c\sqrt{\hat v} + \varepsilon)$, which departs from
  the unscaled value by a factor of $1/(1 + \varepsilon/(c\sqrt{\hat v}))$.
  Making $c$ small therefore makes the epsilon matter, which is the
  measurement in section 4.}
\item A colleague proposes putting the epsilon inside the root
  \enquote{for symmetry}. Give a quantitative argument against it, in units.
  \answerto{Inside, $\varepsilon$ is added to $\hat v$, which has the units of
  a gradient squared. A value of $\num{1e-8}$ is then a floor on $g^{2}$,
  i.e. on $\lvert g \rvert$ of $\num{1e-4}$ \dash{} an enormous gradient to
  treat as noise. Measured, a coordinate with $\lvert g \rvert =
  \val{p20.eps.g}$ loses $\val{p20.eps.short.in}$ per cent of its step where
  the correct placement costs $\val{p20.eps.short.out}$.}
\item Why does the momentum walk in section 2 move \emph{away} from the
  minimum on its first steps, and why does that not contradict the rate?
  \answerto{The average starts at zero and has to be built out of gradients;
  while it is being built its direction lags the current gradient, so it
  carries the walk past the point it was aimed at. The rate is an asymptotic
  statement about the ratio of successive distances once the walk has settled,
  and a ratio says nothing about the constant in front of it. Measured, the
  distance rises to $\val{p20.overshoot}$ times its starting value and the
  tail then falls by $\val{p20.tail.rate}$ per step against a predicted
  $\val{p20.rate.mom}$.}
\item Plain descent's predicted step count was exact at every condition
  number tried and momentum's was not. What is different about the two
  predictions?
  \answerto{Plain descent's iteration has distinct real factors
  $1 - \eta\lambda$ in each eigendirection, so the distance is exactly a
  geometric sequence from the first step and the count follows from the rate
  alone. Momentum's optimal coefficients make the two roots of its two-term
  recurrence coincide, giving a $k\rho^{k}$ term, so the rate is an asymptote
  rather than a description of the whole walk.}
\end{furtherproblems}
